%%==========================================================================
%% This is the very first version of a Streetwise LaTeX template. By
%% chosing the document class, it is easily convertible to a two-column IEEE
%% paper format, provided any conflicting packages and (new)commmands are
%% removed.
%%
%% Erwin de Gelder, January 6, 2018
%%==========================================================================

%%==========================================================================
%% Title:
%% Author(s):
%% To be published in:
%% Date of submission:
%% Date of submission R1:
%% Date of submission R2:
%% Date of final submission:
%%==========================================================================


%% Generic
\documentclass[10pt,final,a4paper,oneside,onecolumn]{article}

%% IEEE
%% Document class options: replace "draftcls" by "final" for final document.
%% All other options may just as well be omitted because they are the default values.
%\documentclass[10pt,final,journal,letterpaper,twoside,twocolumn]{IEEEtran}


%%==========================================================================
%% Document automation
%%==========================================================================

\def\reptitle{Completeness of real-world driving data}
\def\repauthor{Erwin de Gelder, etc.}


%%==========================================================================
%% Packages
%%==========================================================================

\usepackage[a4paper,left=3.5cm,right=3.5cm,top=3cm,bottom=3cm]{geometry} %% change page layout; remove for IEEE paper format
\usepackage[T1]{fontenc}                        %% output font encoding for international characters (e.g., accented)
\usepackage[cmex10]{amsmath}                    %% math typesetting; consider using the [cmex10] option
\usepackage{amssymb}                            %% special (symbol) fonts for math typesetting
\usepackage{amsthm}                             %% theorem styles
\usepackage{dsfont}                             %% double stroke roman fonts: the real numbers R: $\mathds{R}$
\usepackage{mathrsfs}                           %% formal script fonts: the Laplace transform L: $\mathscr{L}$
\usepackage[pdftex]{graphicx}                   %% graphics control; use dvips for TeXify; use pdftex for PDFTeXify
\usepackage{array}                              %% array functionality (array, tabular)
\usepackage{upgreek}                            %% upright Greek letters; add the prefix 'up', e.g. \upphi

\usepackage[utf8]{inputenc}   				 	%% utf8 support (required for biblatex)
\usepackage[style=ieee,doi=false,isbn=false,url=false,date=year,minbibnames=15,maxbibnames=15,backend=biber]{biblatex}
%\renewcommand*{\bibfont}{\footnotesize}		%% Use this for papers
\setlength{\biblabelsep}{\labelsep}
\bibliography{../bib}

\usepackage{stfloats}                           %% improved handling of floats
\usepackage{multirow}                           %% cells spanning multiple rows in tables
%\usepackage{subfigure}                         %% subfigures and corresponding captions (for use with IEEEconf.cls)
%\usepackage{subfig}                             %% subfigures (IEEEtran.cls: set caption=false)
\usepackage{fancyhdr}                           %% page headers and footers
\usepackage[official,left]{eurosym}             %% the euro symbol; command: \euro
\usepackage{appendix}                           %% appendix layout
\usepackage{xspace}                             %% add space after macro depending on context
\usepackage{verbatim}                           %% provides the comment environment
\usepackage[dutch,USenglish]{babel}             %% language support
\usepackage{wrapfig}                            %% wrapping text around figures
\usepackage{longtable}                          %% tables spanning multiple pages
\usepackage{pgfplots}                           %% support for TikZ figures (Matlab)
\usepackage[breaklinks=true,hidelinks,          %% implement hyperlinks (dvips yields minor problems with breaklinks;
            bookmarksnumbered=true]{hyperref}   %% IEEEtran: set bookmarks=false)
%\usepackage[hyphenbreaks]{breakurl}            %% allow line breaks in URLs (don't use with PDFTeX)
\usepackage{lmodern} 
\usepackage{etoolbox}							%% Needed for apptocmd later
\usepackage[capitalize]{cleveref}
\usepackage{units}
\usepackage{subcaption}
\usepackage{csquotes}							%% Quoted texts are typeset according to rules of main language

%%==========================================================================
%% Fancy headers and footers
%%==========================================================================

\newtoggle{standalone}
\toggletrue{standalone}
\pagestyle{fancy}                                       %% set page style
\fancyhf{}                                              %% clear all header & footer fields
\iftoggle{standalone}{%
	\fancyhead[L]{\includegraphics[width=15mm]{Streetwise}} %% define headers (LE: left field/even pages, etc.)
	\fancyhead[R]{\small\emph{\reptitle}}                   %% similar
	\fancyfoot[C]{\thepage}                                 %% define footer
	\setlength{\headheight}{18pt}                           %% increase head height to accommodate an oversized picture in the header
	\renewcommand*{\headrulewidth}{0.25pt}                  %% header line width
	%% Redefine the default "plain" page style (automatically activated by \maketitle, \section, ...)
	\fancypagestyle{plain}{%
		\fancyhf{}
		\fancyhead[C]{\includegraphics[width=30mm]{Streetwise}}
		\renewcommand{\headrulewidth}{0pt}
		\renewcommand{\footrulewidth}{0pt}
		\setlength{\headheight}{36pt}
	}
}{%
	\renewcommand*{\headrulewidth}{0pt}                  	%% No line in this case
	%% Redefine the default "plain" page style (automatically activated by \maketitle, \section, ...)
	\fancypagestyle{plain}{%
		\fancyhf{}
	}
}
\renewcommand*{\footrulewidth}{0pt}                     %% footer line width


%%==========================================================================
%% TikZ figures
%%==========================================================================

\newlength\figurewidth
\setlength\figurewidth{0.35\textwidth}              %% set figure width
\newlength\figureheight
\setlength\figureheight{0.3\textwidth}              %% set figure height
\pgfplotsset{every axis/.append style={
    scaled y ticks=false,
    scaled x ticks=false,
    y tick label style={/pgf/number format/fixed},
    x tick label style={/pgf/number format/fixed},
    legend style={font=\small}},
    compat=1.9}                                     %% PGFPlots package options
\usetikzlibrary{shapes.geometric, arrows, arrows.meta}
\usepgfplotslibrary{groupplots}
%\usetikzlibrary{external}                           %% Create pdf figures from TikZ. Use PDFTeXify ...
%\tikzexternalize[prefix=./tikz/]                    %% ... with --tex-option=--shell-escape switch.
%\tikzset{external/force remake}                    %% force pdf figure update


%%==========================================================================
%% User-defined commands
%%==========================================================================

\newcommand*{\mat}[1]{\mathbf{#1}}                              %% matrix/vector notation
\newcommand*{\matsym}[1]{\boldsymbol{#1}}                       %% matrix/vector notation for Greek letters
\newcommand*{\T}{^{\scriptscriptstyle\mathsf{T}}}               %% transpose operator
\newcommand*{\Hr}{^{\scriptscriptstyle\mathsf{H}}}              %% conjugate transpose operator
\newcommand*{\ud}{\mathrm{\,d}}                                 %% differential operator (upright d)
\newcommand*{\defeq}{\mathrel{\mathop:}=}                       %% definition sign :=
\newcommand*{\eqdef}{=\mathrel{\mathop:}}                       %% definition sign =:
\newcommand*{\ip}[2]{\left\langle#1\,{,}\,#2\right\rangle}      %% inner product
\newcommand*{\real}[1]{\mathrm{Re}(#1)}                         %% real part
\newcommand*{\imag}[1]{\mathrm{Im}(#1)}                         %% imaginary part
\newcommand*{\lsup}[1]{{}^{#1}\!}                               %% left superscript
\newcommand*{\hi}[1]{$^\text{#1}$}                              %% superscript in normal text
\newcommand*{\lo}[1]{$_\text{#1}$}                              %% subscript in normal text
\newcommand*{\w}[1]{\mathrm{#1}}                                %% multiple character super-/subscript in math mode
\newcommand*{\capskip}{\vspace{-12pt}}                          %% caption skip for figures with subfloats
\newcommand*{\etal}{et al.}                                     %% may be required for Natbib bibliography styles
\renewcommand*{\qedsymbol}{$\blacksquare$}                      %% redefine the end-of-proof symbol
\renewcommand*{\labelitemi}{$\bullet$}                          %% first level item list bullet
\renewcommand*{\labelitemii}{$-$}                               %% second level item list bullet
%\renewcommand*{\theenumi}{\textit{\roman{enumi}}}               %% first level enumerator
\renewcommand*{\labelenumi}{\theenumi.}
\renewcommand*{\theenumii}{\textit{\alph{enumii}}}              %% second level enumerator
\renewcommand*{\labelenumii}{\theenumii.}
\DeclareMathOperator{\tr}{tr}                                   %% trace of a matrix
\DeclareMathOperator{\sgn}{sgn}                                 %% signum function
\DeclareMathOperator{\atan}{atan}                               %% arc tangent

\newcommand{\deuclid}{d_{\textup{Euclidian}}}
\newcommand{\ddtw}{d_{\textup{DTW}}}
\newcommand{\imagi}{\textbf{j}}
\newcommand{\profile}[1]{\textbf{v}_{#1}}
\newcommand{\barprofile}[1]{\bar{\textbf{v}}_{#1}}
\newcommand{\parvstart}{v_{\textup{start}}}
\newcommand{\parvend}{v_{\textup{end}}}
\newcommand{\pardv}{v_{\Delta}}
\newcommand{\pardt}{t_{\textup{duration}}}

\crefname{figure}{Figure}{Figures}
\crefname{equation}{}{}
\Crefname{equation}{Equation}{Equations}

%%==========================================================================
%% User-defined environments
%%==========================================================================

\theoremstyle{plain}\newtheorem{definition}{Definition}[section]    %% definition
                    \newtheorem{theorem}{Theorem}[section]          %% theorem
                    \newtheorem{lemma}[theorem]{Lemma}              %% lemma
                    \newtheorem{corollary}[theorem]{Corollary}      %% corollary
                    \newtheorem{assumption}{Assumption}[section]    %% assumption
                    \newtheorem{condition}{Condition}[section]      %% condition
\theoremstyle{definition}\newtheorem{example}{Example}[section]     %% examples
\theoremstyle{remark}\newtheorem{remarkenv}{Remark}[section]        %% remarks
\newenvironment{remark}{\begin{remarkenv}}%
                       {\hfill$\blacklozenge$\end{remarkenv}}       %% end remark with a lozenge
\newenvironment{reviewer}{\itshape}{\upshape}                       %% environment for reviewer's comments


%%==========================================================================
%% Miscellaneous
%%==========================================================================

\graphicspath{{./../}{./figures/}{./more_figures/}} %% (graphicx) directory path for figures
%\setlength{\parindent}{0pt}                        %% no paragraph indentation
%\setlength{\parskip}{2ex}                          %% create empty line between paragraphs
%\interdisplaylinepenalty=2500                      %% (amsmath) allow for page breaks within multiline equations
%\numberwithin{equation}{section}                   %% (amsmath) include section number in equation numbering
%\numberwithin{figure}{section}                     %% (amsmath) include section number in figure numbering
%\numberwithin{table}{section}                      %% (amsmath) include section number in table numbering
\addtolength{\arraycolsep}{-0.5mm}                  %% squeeze matrix columns a little
\fnbelowfloat                                       %% (stfloats) put footnote below a float at the page bottom
\urlstyle{same}                                     %% (hyperref) use current font for URLs
\hypersetup{pdftitle={\reptitle},
            pdfauthor={\repauthor}}                 %% (hyperref) pdf properties title and author
%\raggedbottom                                      %% don't add inter-paragraph spacing to achieve \textheight
%\setlength\subfigcapskip{-3pt}                     %% (subfigure) distance between subfloat and subcaption
%\setlength\subfigbottomskip{-3pt}                  %% (subfigure) distance between subcaption and caption
%\renewcommand*{\subcapsize}{\small}                %% (subfigure) subcaption font size
\renewcommand*{\thesubfigure}{(\alph{subfigure})}   %% (subfig) implement 1(a) instead of 1a ...as subfigure reference
\captionsetup[subfloat]{labelformat=simple}        %% ... as subfigure reference
%\captionsetup[subfloat]{
%    farskip=0pt,
%    nearskip=8pt,
%    captionskip=1.5pt,
%    labelfont={small,bf},
%    textfont=small}                                %% (subfig) subfloat caption format
%\captionsetup[figure]{
%    labelfont={small,bf},
%    textfont=small}                                %% (subfig/caption) figure caption format
%\captionsetup[table]{
%    aboveskip=2pt,
%    labelfont={small,bf},
%    textfont={small,sc}}                           %% (subfig/caption) table caption format
\apptocmd{\thebibliography}{\raggedright}{}{}		%% Suppress badness warnings in bibliography


% Table stuff
\usepackage{booktabs}
\usepackage{tabularx}
\setlength{\heavyrulewidth}{0.1em}
\newcommand{\otoprule}{\midrule[\heavyrulewidth]}

%%==========================================================================
%% Begin document
%%==========================================================================

\begin{document}

\selectlanguage{USenglish}

\title{\textbf{\reptitle}}
\author{\repauthor}
\date{\today}
\maketitle

\tableofcontents

\newpage

\section{Introduction}
\label{sec:introduction}

% Scenario-based approach for assessment of AVs -> this scenarios are important
An important aspect in the development of automated vehicles (AVs) is the assessment of quality and performance aspects of the AVs, such as safety, comfort, and efficiency \cite{bengler2014threedecades, stellet2015taxonomy, wachenfeld2016release, putz2017pegasus, roesener2016scenariobased, kompass2015sicherheitsveranderung}. 
For legal and public acceptance, it is important that there is a clear definition of system performance and that there are quantitative measures for the system quality. 
The more traditional methods \cite{response2006code, ISO26262}, used for evaluation of driver assistance systems, are no longer valid for the assessment of quality and performance aspects of an AV \cite{wachenfeld2016release}. 
Therefore, a scenario-based approach is proposed \cite{roesener2016scenariobased, putz2017pegasus, kompass2015sicherheitsveranderung}. 
For the scenario-based assessment, proper specification of scenarios is crucial since they are directly reflected in test cases used for scenario-based assessment \cite{stellet2015taxonomy}.

% Completeness
In order to draw conclusions on how an automated vehicle would perform in real-world traffic, it is necessary to know how representative the scenario database, which is used for the scenario-based assessment of the automated vehicle, is. Therefore it is important to quantify how complete the scenario database is \cite{geyer2014, alvarez2017prospective, stellet2015taxonomy}.

% Purpose of this document
This document describes different methods for quantifying the completeness of a dataset. To illustrate the different methods, use is made of artificial data. The advantage of using artificial data is that it is easy to obtain a huge amount of data. The data consists of the speed of a vehicle sampled at $\unit[10]{Hz}$. Furthermore, the acceleration is estimated by filtering the derivative of the speed. A portion of the data is shown in \cref{fig:artificial data}. Note that the data is not very realistic, because in real-life the relative time a vehicle is cruising, i.e., driving at a constant speed, will be higher. However, for the purpose of illustrating the different methods, the artificial data suffices.

\setlength\figurewidth{0.8\linewidth}              %% set figure width
\setlength\figureheight{0.3\textwidth}              %% set figure height
\begin{figure}[b]
	\centering
	% This file was created by matplotlib2tikz v0.6.14.
\begin{tikzpicture}

\definecolor{color0}{rgb}{0.12156862745098,0.466666666666667,0.705882352941177}

\begin{axis}[
xlabel={Time [s]},
ylabel={Speed [m/s]},
xmin=0, xmax=500,
ymin=-1.3007710998959, ymax=27.3161930978138,
width=\figurewidth,
height=\figureheight,
tick align=outside,
tick pos=left,
xmajorgrids,
x grid style={white!69.019607843137251!black},
ymajorgrids,
y grid style={white!69.019607843137251!black}
]
\addplot [line width=2.0pt, color0, forget plot]
table {%
0 0.02495226681389
1 0.41412915559112
2 1.20415249714665
3 2.04559320305136
4 2.96349908784175
5 4.05774866180894
6 5.12866058530594
7 6.33534420532051
8 7.02440786605288
9 7.25058904747158
10 7.35198356916793
11 7.53558457779503
12 7.53558457779503
13 7.53558457779503
14 7.53558457779503
15 7.53558457779503
16 7.53558457779503
17 7.53558457779503
18 7.53558457779503
19 7.53558457779503
20 7.53558457779503
21 7.53558457779503
22 7.53558457779503
23 7.53558457779503
24 7.53558457779503
25 7.53558457779503
26 7.52549984694471
27 7.42982944658301
28 7.37313876723882
29 7.35259530108769
30 7.19363515011014
31 7.12198980592494
32 6.92134867417635
33 6.74889592880281
34 6.54825446130259
35 6.34655984429617
36 6.1292326420585
37 5.94098500435689
38 5.69494203639759
39 5.47554865744759
40 5.26543085890218
41 5.02274943310738
42 4.8156281044786
43 4.57639024257758
44 4.28401370382455
45 4.07446246924165
46 3.83874816162199
47 3.52668706196945
48 3.33329751272105
49 3.08057217745097
50 2.80018403051738
51 2.63291621487606
52 2.44146863581963
53 2.26994966908652
54 2.09016527687301
55 1.91280740001334
56 1.71917366974441
57 1.51941733687503
58 1.32376222611795
59 1.10465280538669
60 0.923368041502252
61 0.762297859725797
62 0.549755191500437
63 0.392930330836383
64 0.274226520723161
65 0.184824745227919
66 0.149237927130451
67 0.0642741344089064
68 0
69 0
70 0
71 0
72 0
73 0
74 0
75 0
76 0
77 0
78 0
79 0
80 0
81 0
82 0
83 0
84 0
85 0
86 0
87 0
88 0
89 0
90 0
91 0
92 0
93 0
94 0
95 0
96 0
97 0
98 0
99 0
100 0
101 0
102 0
103 0.0235510472071443
104 0.588314158964943
105 1.19484288421782
106 1.67410401366337
107 2.06507202894072
108 2.48864779705037
109 2.91059763466345
110 3.34263982019536
111 4.12262318863332
112 5.51238060675081
113 7.07373578415886
114 8.77992636596677
115 10.7249260776893
116 10.4672392154854
117 10.0757932891272
118 9.56215224844886
119 9.14376966264727
120 8.67920735721271
121 8.25774714992156
122 7.84226331300954
123 7.33128031296112
124 6.93576950755434
125 6.68855923520865
126 6.64314153361708
127 6.64314153361708
128 6.64314153361708
129 6.64314153361708
130 6.64314153361708
131 6.64314153361708
132 6.64314153361708
133 6.64314153361708
134 6.64314153361708
135 6.64314153361708
136 6.64314153361708
137 6.64314153361708
138 6.64314153361708
139 6.64314153361708
140 6.77074145955114
141 6.89073386191043
142 7.06366690065455
143 7.18749913095217
144 7.32973394535564
145 7.4975372581904
146 8.01658524200401
147 8.70663554583158
148 9.08290544592541
149 9.28402078519807
150 9.46806007221825
151 9.7223105154507
152 10.1818934142756
153 10.7413384655312
154 11.2984796098625
155 11.9176586938559
156 12.5194171938484
157 13.1673989001897
158 13.8185634610461
159 14.4960619237495
160 15.1351324964136
161 15.9025494632512
162 16.7470706337469
163 17.7725443330002
164 19.0158267752213
165 20.1684419727177
166 21.1562915795093
167 22.0079639517925
168 22.8376271644745
169 24.0096642317662
170 23.4896155489254
171 22.5931385655683
172 21.5635094835419
173 20.5878943722396
174 19.6332161485003
175 19.3518373981516
176 19.3518373981516
177 19.3518373981516
178 19.3518373981516
179 19.3518373981516
180 19.3518373981516
181 19.3518373981516
182 19.3518373981516
183 19.3518373981516
184 19.3518373981516
185 19.3518373981516
186 19.3518373981516
187 19.3518373981516
188 19.4025552783945
189 20.2546208335804
190 21.7480743117486
191 23.5130673481488
192 24.896501678608
193 25.5605191776938
194 26.0154219979179
195 26.0154219979179
196 26.0154219979179
197 26.0154219979179
198 26.0154219979179
199 26.0154219979179
200 26.0154219979179
201 26.0154219979179
202 26.0154219979179
203 26.0154219979179
204 26.0154219979179
205 26.0154219979179
206 26.0154219979179
207 26.0154219979179
208 25.7111559691589
209 25.0607980681135
210 24.4066165406515
211 23.6002186510014
212 22.7977079879777
213 22.0465377827481
214 21.5102868770601
215 20.8941465889376
216 20.4073136737359
217 19.9737336244629
218 19.7531383622051
219 19.6295287448304
220 19.6295287448304
221 19.6295287448304
222 19.6295287448304
223 19.6295287448304
224 19.6295287448304
225 19.6295287448304
226 19.6295287448304
227 19.6295287448304
228 19.6295287448304
229 19.6295287448304
230 19.6295287448304
231 19.5715018533823
232 19.0468180577465
233 18.5350770416726
234 18.0998583315973
235 17.7876979316757
236 17.3914555497333
237 17.1015741384041
238 16.7190056305012
239 16.3919276175824
240 16.1888654412743
241 16.1888654412743
242 16.1888654412743
243 16.1888654412743
244 16.1888654412743
245 16.1888654412743
246 16.1888654412743
247 16.1888654412743
248 16.1888654412743
249 16.1888654412743
250 16.1888654412743
251 16.1888654412743
252 16.1888654412743
253 16.1888654412743
254 16.1888654412743
255 16.1888654412743
256 16.1888654412743
257 16.244759975667
258 16.8376672865179
259 17.5051018539717
260 18.2576967555946
261 18.9737845229909
262 19.0598339898966
263 19.3737393645219
264 19.7780051732413
265 19.7750646496584
266 19.5818271855058
267 19.528532939132
268 19.3897714092076
269 19.1730484775524
270 19.0513159615543
271 18.8363440657725
272 18.6405685667133
273 18.4525606006436
274 18.1930288389244
275 18.0740734911633
276 17.7522965652754
277 17.5547049457927
278 17.3370619938893
279 16.9285281165711
280 16.7876828455576
281 16.4786935853557
282 16.2479875656713
283 15.9129409464543
284 15.5947025134267
285 15.3516940963531
286 14.9490347618845
287 14.7103429485559
288 14.3484474459726
289 14.0220341010075
290 13.7335911333508
291 13.3242325132686
292 13.0725891327197
293 12.7156900135848
294 12.3962597387869
295 12.0950710845955
296 11.8470410830888
297 11.4702597792302
298 11.1122290580795
299 10.8289092166501
300 10.4307816182326
301 10.231867081881
302 9.79736215562572
303 9.48660125092045
304 9.18812089915069
305 8.80291457171979
306 8.53319665952585
307 8.18008216362833
308 7.88378830262853
309 7.53092592243636
310 7.18615042641893
311 6.91527427402575
312 6.53860577725202
313 6.28159780446968
314 5.91660645755556
315 5.5632972742191
316 5.30163674433125
317 4.97631348956604
318 4.73939539813799
319 4.43207675746579
320 4.13442934579078
321 3.92553464856429
322 3.63663116278089
323 3.47398684094807
324 3.20007407439704
325 2.96601681879235
326 2.7549204602362
327 2.48956742290133
328 2.31407933796254
329 2.05084871091424
330 1.8603454909734
331 1.63068983844745
332 1.47089449479967
333 1.1962613977489
334 1.00149028199273
335 0.778947418938918
336 0
337 0
338 0
339 0
340 0
341 0
342 0
343 0
344 0
345 0
346 0
347 0
348 0
349 0
350 0
351 0
352 0
353 0
354 0
355 0
356 0
357 0
358 0
359 0
360 0
361 0
362 0
363 0
364 0
365 0
366 0
367 0
368 0
369 0
370 0
371 0
372 0
373 0
374 0
375 0
376 0
377 0
378 0
379 0
380 0
381 0
382 0
383 0
384 0
385 0
386 0
387 0.00892171360271377
388 0.0981386339096986
389 0.362222895776485
390 0.806968384180307
391 1.33290797599514
392 1.71118099526973
393 1.98640461852989
394 2.21926758459486
395 2.37940502248795
396 2.79428684471596
397 2.8749967233432
398 2.89999941338979
399 3.03917831746172
400 3.34429269924698
401 3.61685873113329
402 3.82337750933905
403 4.09145640525861
404 4.3662212620465
405 4.87253283260646
406 5.48946755359851
407 6.01672383070859
408 6.55246704472216
409 7.18725031364922
410 7.81402483661992
411 8.4644207015304
412 9.03584991060398
413 9.57873477076918
414 9.99359583305322
415 10.4066656896986
416 10.7448022582683
417 11.0307436864592
418 10.903825648276
419 10.1346057409828
420 9.1460041169704
421 8.16759985745562
422 7.51621256146151
423 7.51621256146151
424 7.51621256146151
425 7.51621256146151
426 7.51621256146151
427 7.51621256146151
428 7.51621256146151
429 7.51621256146151
430 7.51621256146151
431 7.51621256146151
432 7.51621256146151
433 7.51621256146151
434 7.51621256146151
435 7.41353993523616
436 6.66263204457012
437 5.75267170116231
438 4.86183333207651
439 3.91664049713349
440 3.00768896969473
441 2.03631660135979
442 1.02672437477843
443 0
444 0
445 0
446 0
447 0
448 0
449 0
450 0
451 0
452 0
453 0
454 0
455 0
456 0
457 0
458 0.478592642672584
459 0.770708092597168
460 1.03872527629063
461 1.63218507114834
462 2.36813306141615
463 3.09204464323763
464 3.52951065764046
465 3.55574416565908
466 3.55574416565908
467 3.55574416565908
468 3.55574416565908
469 3.53129014676358
470 3.19387417395677
471 2.77685196794471
472 2.44515921876903
473 2.01200601908199
474 1.61865749700049
475 1.23571423657206
476 0.858497821150841
477 0.50495344521193
478 0.208380993476534
479 0.036681390349861
480 0
481 0
482 0
483 0
484 0
485 0
486 0
487 0
488 0
489 0
490 0
491 0
492 0
493 0
494 0
495 0
496 0
497 0
498 0
499 0
};
\end{axis}

\end{tikzpicture}
	\caption{Artificial data of the speed of a vehicle.}
	\label{fig:artificial data}
\end{figure}

% Structure of document


\section{Sample-based methods}
\label{sec:sample based}


\subsection{Comparing probability density functions}
\label{sec:sample based pdf}

According to Wang \etal\ \cite{wang2017much}, the question of ``how much driving data is sufficient to address problems such as the cause of accidents, distraction and inattention, eco-driving styles, modeling driver behavior, and the effects of driver assistance systems on driver behavior'' has not been fully studied. Therefore, Wang \etal\ \cite{wang2017much} present a method for quantifying the completeness regarding naturalistic driving data. 

Wang \etal\ \cite{wang2017much} use a statistical method for determining the appropriate amount of data. Based on the data, a probability density function (PDF) is estimated. The point at which the PDF converges determines the amount of appropriate data. 

Let the PDF of a certain state $x \in \mathds{R}^D$ be denoted by $f(x)$. Based on the $n$ data points $\textbf{x}=\{x_i\}_{i=1}^n$, the PDF $f(x)$ can be estimated. This estimation is denoted by $\hat{f}(x;n)$. For estimating $f(x)$, Kernel Density Estimation (KDE) \cite{rosenblatt1956remarks, parzen1962estimation} is employed, i.e.,
\begin{equation}
	\hat{f}(x;n) = \sum_{i=1}^n \frac{1}{h^D n} K \left( \frac{x - x_i}{h} \right),
\end{equation}
where $h$ denotes the so-called bandwidth and $K(\cdot)$ represents the kernel. A Gaussian kernel is used, i.e.,
\begin{equation}
	K(x) = \frac{1}{\sqrt{2\pi}} \exp \left\{ -\frac{||x||^2}{2}\right\}.
\end{equation}
For the bandwidth, the Silverman rule is used, i.e. $h=1.06\cdot\hat{\sigma}\cdot n^{-1/5}$ \cite{silverman1986density}, where $\hat{\sigma}$ is the standard deviation of the data $\{x_i\}_{i=1}^n$.

\begin{remark}
	The power of the Silverman rule is in its simplicity. The disadvantage, however, is that it tends to oversmooth the density when it is used for multimodal and highly skewed densities \cite{silverman1986density}. When comparing two estimated PDFs, they will be more similar when both are oversmoothed. As the comparison of two estimated PDFs leads to the quantification of completeness, the use of the Silverman rule might lead to an overestimation of the degree of completeness.
\end{remark}

The distribution $f(x;n)$ will change as $n$ increases. It is expected, however, that the change will decrease as $n$ increases. Therefore, the distribution $\hat{f}(x;n+m)$ with $m \in \mathds{R}^+$ is compared with $\hat{f}(x;n)$. For the comparison of two PDFs, the Kullback-Leibler divergence is employed \cite{kullback1951, bishop2006pattern}. For the two distributions $\hat{f}(x;n+m)$ and $\hat{f}(x;n)$, the Kullback-Leibler divergence is as follows:
\begin{equation}
	KL\left( \hat{f}(x;n+m) || \hat{f}(x;n) \right) = \int \hat{f}(x;n+m) \ln \left( \frac{\hat{f}(x;n+m)}{\hat{f}(x;n)} \right) \ud x.
\end{equation}

The Kullback-Leibler divergence is always positive and zero if and only if $\hat{f}(x;n+m) = \hat{f}(x;n)$. Wang~\etal\ determine the proper amount of driving data so that $KL(\hat{f}(x;n+m)||\hat{f}(x;n)$ changes very slightly if more data samples are to be added, i.e.,
\begin{equation} \label{eq:completeness wang}
	\left| KL\left(\hat{f}(x;n+m)||\hat{f}(x;n)\right) - KL\left(\hat{f}(x;n+2m)||\hat{f}(x;n+m)\right) \right| \leq \epsilon,
\end{equation}
for a small value of $\epsilon\in\mathds{R}^+$. The proper number of data-points $n$ equals the minimum value of $n$ for which \cref{eq:completeness wang} is satisfied.

\begin{remark}
	Although Wang~\etal\ use \cref{eq:completeness wang} to determine the appropriate number of data points $n$, it seems a bad choice as \cref{eq:completeness wang} is satisfied if the difference between $\hat{f}(x;n+m)$ and $\hat{f}(x;n)$ is almost similar to the difference between $\hat{f}(x;n+m)$ and $\hat{f}(x;n+2m)$. A better should would be to use
	\begin{equation}
		KL\left(\hat{f}(x;n+m) || \hat{f}(x;n) \right) \leq \epsilon,
	\end{equation}
	for a small value of $\epsilon\in\mathds{R}^+$.
\end{remark}

\begin{example}
	\Cref{eq:completeness wang} is applied for the artificial data as shown in \cref{fig:artificial data}. When using similar values for $m$ and $\epsilon$ as in \cite{wang2017much}, i.e., $m=2000$ and $\epsilon=10^{-4}$, and the velocity as the feature $x$, it appears that 27 minutes of data suffices. 
	
	The result depends on the choice of $m$ and on the feature that is used. \Cref{tab:wang result} lists the results for different values of $m$, where both the speed and the acceleration are considered as features. Considering the estimations regarding the different features, Wang \etal\ use the maximum as final estimation of the appropriate amount of data. The maximum required time is shown in the last column of \cref{tab:wang result}. \Cref{tab:wang result} shows that the final estimated of the required data heavily depends on $m$.
		
	\begin{table}
		\centering
		\caption{Appropriate data according to \cref{eq:completeness wang}.}
		\label{tab:wang result}
		\begin{tabular}{rrrr}
			\toprule
			& \multicolumn{2}{c}{Feature} & \\
			$m$ & Speed & Acceleration & Maximum \\ \otoprule
			500 & $\unit[7]{min}$ & $\unit[5]{min}$ & $\unit[7]{min}$ \\
			100 & $\unit[7]{min}$ & $\unit[12]{min}$ & $\unit[12]{min}$ \\
			2000 & $\unit[27]{min}$ & $\unit[23]{min}$ & $\unit[27]{min}$ \\
			4000 & $\unit[67]{min}$ & $\unit[27]{min}$ & $\unit[67]{min}$ \\
			6000 & $\unit[50]{min}$ & $\unit[40]{min}$ & $\unit[50]{min}$ \\
			8000 & $\unit[93]{min}$ & $\unit[40]{min}$ & $\unit[93]{min}$ \\
			10000 & $\unit[100]{min}$ & $\unit[33]{min}$ & $\unit[100]{min}$ \\
			\bottomrule
		\end{tabular}
	\end{table}
\end{example}

\subsection{Likelihood of new samples}

This method is similar to the previous method in that a PDF is estimated using $n$ samples, i.e., $\hat{f}(x;n)$. But instead of comparing $\hat{f}(x;n)$ with $\hat{f}(x;n+m)$, the $m$ new samples are directly compared with first $n$ samples by determining the likelihood of the new $m$ samples based on the estimated PDF $\hat{f}(x;n)$. 

Assuming that the samples are independent, the likelihood is simply the product of $\hat{f}(x_i;n)$ with $i=\{n+1, \ldots, n+m\}$. To get a value that is less dependent on the choice of $m$, the following value is used:
\begin{equation}
	J(n) = \left( \prod_{i=n+1}^{n+m} \hat{f}(x_i;n) \right)^{\frac{1}{m}}.
\end{equation}

\begin{remark}
	Obviously, consecutive samples are not independent.
\end{remark}

\printbibliography

\end{document} 